\chapter{Considerações Finais}\label{cap:Considerações}

\section{Contribuição e alcance}

Esta tese combina densidade espacial (PDFe), EOF univariado e MEOF para quantificar, por fase do ciclo de vida e região de gênese, a separação espacial entre vento e precipitação nos ciclones extratropicais do Atlântico Sul-Ocidental. A constatação central, de que a intensidade dinâmica do vórtice deixa de predizer linearmente a magnitude e a localização de seus impactos de superfície no subconjunto p90, e seu respaldo estatístico e geográfico, já foram desenvolvidos no Capítulo~\ref{ch:conclusiones}. A confiança de que esses padrões espaciais não são artefatos da amostragem se sustenta na reamostragem Bootstrap-t (Seção~\ref{subsec:bootstrap_significancia}), que documenta áreas estatisticamente significativas no EOF1 tanto na População Global quanto no subconjunto p90, estendendo-se neste último também ao EOF2 e EOF3, os mesmos modos que capturam a reorganização de mesoescala nos sistemas intensos, em alguns casos com maior área significativa que o próprio EOF1, o que respalda que o sinal relatado é físico e recorrente, e não ruído estocástico. As limitações e linhas futuras que seguem delimitam o alcance do que foi apresentado neste capítulo.

\section{Limitações}

Os resultados estão sujeitos à subestimação do ERA5 em extremos ($-10.2\%$ em vento e $-22.6\%$ em precipitação para os ciclones mais intensos), de modo que as magnitudes absolutas do subconjunto p90 são estimativas conservadoras. As estruturas de mesoescala responsáveis pelos máximos de vento em ciclones severos, como a CCB, os \textit{sting jets} ou as descargas descendentes, operam em escalas menores que a resolução do ERA5 ($\sim$31~km), limitando a resolução de sua geometria interna nos protótipos. A padronização do domínio centrado no ciclone em $\pm 9{,}5^{\circ}$ pode truncar a estrutura em arco de precipitação nos ciclones de maior escala horizontal, subestimando a extensão real do campo pluviométrico periférico.

O tamanho amostral do subconjunto p90 difere substancialmente entre regiões (198 ciclones em ARG, frente a 64 em SBR e 67 em LPB), de modo que as densidades periféricas das PDFe e os autovetores do EOF em SBR e LPB estão sujeitos a maior incerteza de amostragem que os de ARG. A avaliação de significância estatística mediante Bootstrap-t se concentrou na fase de maturidade, o momento de maior interesse para o argumento central da tese, de modo que os padrões EOF das fases restantes são apresentados como caracterização descritiva sem a mesma validação formal, e o modo multivariado (MEOF) se apoia na significância já estabelecida de seus componentes univariados na maturidade sem um teste multivariado independente. O próprio EOF, por ser uma decomposição linear, poderia diluir em um único modo diferenças que na realidade correspondam a subpopulações de ciclones estruturalmente discretas, embora o aplanamento do espectro de autovalores observado nos sistemas p90 sugira que o método está capturando, e não ocultando, essa complexidade.

A análise se circunscreve aos campos de superfície. A extensão à coluna tridimensional permitiria vincular as estruturas identificadas com seus mecanismos geradores.

O catálogo de trajetórias empregado \citep{gramcianinov2020analysis} constitui uma única fonte de detecção e rastreamento entre vários esquemas objetivos disponíveis para o Atlântico Sul (Seção~\ref{sec:tb_vorticidad}); embora baseado em vorticidade relativa a 850~hPa, um critério amplamente validado, não se pode descartar que outro catálogo modifique a composição exata das populações PG e p90. Espera-se, no entanto, que os sinais estruturais aqui documentados (a separação em quadrantes, a fragmentação do espectro de autovalores em p90) sejam robustos frente a essa escolha, na medida em que dependem da física do ciclone e não do algoritmo de detecção particular.

Diferentemente de estudos do Hemisfério Norte que rotacionam cada composite para alinhar a direção de deslocamento do sistema, por exemplo em direção ao leste, antes de calcular a média \citep{rudeva2011composite, priestley2022improved}, o referencial centrado no sistema aqui empregado mantém a orientação geográfica original (Seção~\ref{subsec:tb_lagrangiano}). Essa escolha preserva a interpretação direta dos quadrantes em termos cardinais (noroeste, sudeste), mas não corrige a variabilidade na direção de translação dos ciclones, o que poderia diluir, nunca exagerar artificialmente, as assimetrias estruturais relatadas; que estas emerjam com clareza apesar dessa fonte adicional de dispersão respalda a robustez do sinal detectado.

Adicionalmente, a exigência de um ciclo de vida completo com as quatro fases e sem reintensificações (Seção~\ref{subsec:grupo_global}) exclui sistemas de evolução complexa, truncada ou de vida curta, que poderiam representar uma fração não desprezível dos ETC reais da região. Os resultados aqui obtidos caracterizam, portanto, o subconjunto de ciclones com evolução estruturalmente completa, não a totalidade da atividade ciclônica do Atlântico Sul-Ocidental.

\section{Perspectivas futuras}

Como linhas futuras, seria possível validar a geometria dos protótipos com observações de alta resolução (boias, estações costeiras, satélites) para ajustar os vieses do ERA5 no subconjunto p90; incorporar vento vetorial em 850~hPa, temperatura potencial equivalente e vorticidade potencial à análise MEOF permitiria relacionar os padrões de superfície com sua evolução em altura, incluindo as correntes transportadoras (WCB, CCB) e o \textit{sting jet}; os protótipos também serviriam para definir limiares combinados de vento e precipitação e estimar a probabilidade de impactos conjuntos nas costas do Brasil, Uruguai e Argentina, distinguindo o perigo do vento (noroeste) do da chuva (sudeste-sul) segundo a posição do ciclone em relação à costa; simulações de alta resolução com diferentes cenários climáticos permitiriam avaliar se o aquecimento global modifica a separação de quadrantes aqui documentada; e estender a reamostragem Bootstrap-t a todas as fases do ciclo de vida e ao MEOF, junto com um registro mais longo ou dados de modelos de alta resolução para SBR e LPB, reduziria a incerteza de amostragem frente a ARG.
